%%%PAGE 284 rot=0 legibility=good summary=粘贴的20分综合题（区域I水平电场、区域II竖直电场加磁场的交替运动），含手绘场区图和旁注答案；页面下半空白
%%%BLOCK id=284-01 ch=11 sec=复合场·电场磁场交替区域 kind=problem alt=09,04 cont=none y=0.00-0.40
% 题干为粘贴的印刷题；手写对“E1”“E2”“匀强磁场”“正电”“初速度v0竖直向上抛出”“(-L,0)”画了方框/圈，对“区域II中做匀速圆周运动”“恰好通过坐标原点O回到区域I”“交替在区域I和区域II中运动”“第二次进入区域II时的位置坐标”“轨迹半径r2”画了下划线
% 右下角折页露出背面一小幅手绘图（斜线与一条末端带圈的曲线），属另一面内容，未转录
\begin{problem}[（20分）]
如图所示，在竖直平面内建立直角坐标系 $xOy$，$x$ 轴上方的区域I内存在水平向右的匀强电场 $E_1$，$x$ 轴下方的区域II内存在竖直向上的匀强电场 $E_2$ 和垂直于竖直平面向外的匀强磁场。将一个带正电的小球从 $x$ 轴上坐标为 $(-L,0)$ 的 $P$ 点以初速度 $v_0$ 竖直向上抛出，小球垂直通过 $y$ 轴进入第一象限，然后通过 $x$ 轴进入区域II中做匀速圆周运动，恰好通过坐标原点 $O$ 回到区域I，之后小球可以交替在区域I和区域II中运动。已知小球第三次进入区域II时的动能为其初动能的4倍，重力加速度未知。求：

（1）$E_1$ 和 $E_2$ 的大小之比；\quad （手写）\unclear{$\sqrt{3}:6$}% CHECK: 手写为“√3 ∶ 6”，6 的起笔很长，也可能是“√3 ∶ √6”

（2）小球第一次在区域II中运动的时间 $t_1$；\quad （手写）$\dfrac{2\pi L}{v_0}$

（3）小球第二次进入区域II时的位置坐标和第二次在区域II中做匀速圆周运动的轨迹半径 $r_2$

（手写）$r_2=\dfrac{\sqrt{21}}{2}L$

%FIG p284_a 0.645 0.015 1.0 0.235
\notefig[0.35]{figures/p284_a.png}
\end{problem}
%%%END
%%%PAGEEND 284
